% GENERATED FILE. Edit canonical lessons, then run npm run generate:books.
% canonical-source-hash: fnv1a64:7bd196e4ddc6d042
% canonical-lessons: ZH-C268-baoming, ZH-C268-koushi, ZH-C268-jige, ZH-C268-daan, ZH-C268-jieshi

\chapter{Exams}
\label{ch:zh-exams}
\hlchaptermodality{268}

\begin{chapteropening}
\textbf{By the end of this chapter:} \emph{I can sign up for an exam, talk about an oral exam, passing and the answers, and ask for an explanation.}

Complete the last lesson of chapter 268: I can sign up for an exam, talk about an oral exam, passing and the answers, and ask for an explanation.
\end{chapteropening}

\section[bàomíng]{bàomíng — to sign up; to enter (one's name)}

\label{lesson:ZH-C268-baoming}



\begin{quote}
Before the new one: say the Mandarin for a joke, then the Mandarin for funny.
\end{quote}

\subsection*{You'll want to know: bàomíng}
\emph{bàomíng} — \textquotedblleft{}to sign up; to enter (one's name)\textquotedblright{}.

\begin{quote}
\emph{nǐ bàomíng le ma} — Have you signed up?
\end{quote}

Its characters come later; for now it is a word for the ear.

\begin{cousinweb}
\emph{bàomíng} is built from \emph{bào} 'to report' and \zh{名} 'a name', both inherited from Middle Chinese, \emph{pawH} and \emph{mjieng} (Guangyun, in Baxter's transcription), which glosses the first 'to report'.
\end{cousinweb}

\begin{grammarlens}[title={What you've built}]
One more word for saying what I do.
\end{grammarlens}

\subsection*{Your turn}
\begin{itemize}
\raggedright
  \item \emph{Say it:} \emph{bàomíng}
  \item \emph{Say it:} \emph{bàomíng}, once more
  \item \emph{Say it:} \emph{nǐ bàomíng le ma} — Have you signed up?
  \item \emph{Recall:} say the Mandarin for a joke, then the Mandarin for funny, then say \emph{bàomíng} again
\end{itemize}

\begin{tcolorbox}[breakable,colback=teal!4,colframe=teal!35!black,title={Before you move on}]
What does bàomíng mean? (\textquotedblleft{}To sign up; to enter (one's name)\textquotedblright{}.) Say it once more.
\end{tcolorbox}
\section[\texorpdfstring{\zh{口试} (kǒushì)}{口试 (kǒushì)}]{\zh{口试} (kǒushì) — oral exam}

\label{lesson:ZH-C268-koushi}



\begin{quote}
Before the new one: say the Mandarin for funny, then the Mandarin for to sign up.
\end{quote}

\subsection*{You'll want to know: \zh{口试}}
\textbf{\zh{口试}} — \emph{kǒushì} — \textquotedblleft{}oral exam\textquotedblright{}.

\begin{quote}
\textbf{\zh{明天我们学校有口试}} \emph{míngtiān wǒmen xuéxiào yǒu kǒushì} — Our school has oral exams tomorrow.
\end{quote}

\begin{cousinweb}
\textbf{\zh{口试}} (\emph{kǒushì}) is built from \zh{口} 'the mouth' and \zh{试} 'a test', both inherited from Middle Chinese, \emph{khuwX} and \emph{syiH} (Guangyun, in Baxter's transcription), which glosses the second 'to use'. The sense 'a test' is not given there.
\end{cousinweb}

\begin{grammarlens}[title={What you've built}]
One more word for saying what I do.
\end{grammarlens}

\subsection*{Your turn}
\begin{itemize}
\raggedright
  \item \emph{Say it:} \emph{kǒushì}
  \item \emph{Say it:} \emph{kǒushì}, once more
  \item \emph{Say it:} \emph{míngtiān wǒmen xuéxiào yǒu kǒushì} — Our school has oral exams tomorrow.
  \item \emph{Recall:} say the Mandarin for funny, then the Mandarin for to sign up, then say \emph{kǒushì} again
\end{itemize}

\begin{tcolorbox}[breakable,colback=teal!4,colframe=teal!35!black,title={Before you move on}]
What does \zh{口试} mean? (\textquotedblleft{}Oral exam\textquotedblright{}.) Say it once more.
\end{tcolorbox}
\section[jígé]{jígé — to pass an exam}

\label{lesson:ZH-C268-jige}



\begin{quote}
Before the new one: say the Mandarin for to sign up, then the Mandarin for an oral exam.
\end{quote}

\subsection*{You'll want to know: jígé}
\emph{jígé} — \textquotedblleft{}to pass an exam\textquotedblright{}.

\begin{quote}
\emph{tā kǎoshì méi jígé} — He didn't pass the exam.
\end{quote}

Its characters come later; for now it is a word for the ear.

\begin{cousinweb}
\emph{jígé} is built from \emph{jí} 'to reach' and \emph{gé} 'a standard', both inherited from Middle Chinese, \emph{gip} and \emph{kaek} (Guangyun, in Baxter's transcription), which glosses the first 'to reach' and the second 'a model' and 'a measure'.
\end{cousinweb}

\begin{grammarlens}[title={What you've built}]
One more word for saying what I do.
\end{grammarlens}

\subsection*{Your turn}
\begin{itemize}
\raggedright
  \item \emph{Say it:} \emph{jígé}
  \item \emph{Say it:} \emph{jígé}, once more
  \item \emph{Say it:} \emph{tā kǎoshì méi jígé} — He didn't pass the exam.
  \item \emph{Recall:} say the Mandarin for to sign up, then the Mandarin for an oral exam, then say \emph{jígé} again
\end{itemize}

\begin{tcolorbox}[breakable,colback=teal!4,colframe=teal!35!black,title={Before you move on}]
What does jígé mean? (\textquotedblleft{}To pass an exam\textquotedblright{}.) Say it once more.
\end{tcolorbox}
\section[dá'àn]{dá'àn — answer to a question}

\label{lesson:ZH-C268-daan}



\begin{quote}
Before the new one: say the Mandarin for an oral exam, then the Mandarin for to pass an exam.
\end{quote}

\subsection*{You'll want to know: dá'àn}
\emph{dá'àn} — \textquotedblleft{}answer to a question\textquotedblright{}.

\begin{quote}
\emph{nǐ de dá'àn duì le} — Your answer is right.
\end{quote}

Its characters come later; for now it is a word for the ear.

\begin{cousinweb}
\emph{dá'àn} is built from \emph{dá} 'to answer' and \emph{àn}, both inherited from Middle Chinese, \emph{top} and \emph{'anH} (Guangyun, in Baxter's transcription), which glosses the first 'to match' and the second 'a low table'. How the pair came to mean an answer is not settled.
\end{cousinweb}

\begin{grammarlens}[title={What you've built}]
One more word for saying what I do.
\end{grammarlens}

\subsection*{Your turn}
\begin{itemize}
\raggedright
  \item \emph{Say it:} \emph{dá'àn}
  \item \emph{Say it:} \emph{dá'àn}, once more
  \item \emph{Say it:} \emph{nǐ de dá'àn duì le} — Your answer is right.
  \item \emph{Recall:} say the Mandarin for an oral exam, then the Mandarin for to pass an exam, then say \emph{dá'àn} again
\end{itemize}

\begin{tcolorbox}[breakable,colback=teal!4,colframe=teal!35!black,title={Before you move on}]
What does dá'àn mean? (\textquotedblleft{}Answer to a question\textquotedblright{}.) Say it once more.
\end{tcolorbox}
\section[jiěshì]{jiěshì — to explain}

\label{lesson:ZH-C268-jieshi}



\begin{quote}
Before the new one: say the Mandarin for to pass an exam, then the Mandarin for the answer to a question.
\end{quote}

\subsection*{You'll want to know: jiěshì}
\emph{jiěshì} — \textquotedblleft{}to explain\textquotedblright{}.

\begin{quote}
\emph{zhège dá'àn nǐ néng jiěshì ma} — Can you explain this answer?
\end{quote}

Its characters come later; for now it is a word for the ear.

\begin{cousinweb}
\emph{jiěshì} is built from \emph{jiě} and \emph{shì}, both inherited from Middle Chinese, \emph{keaX} and \emph{syek} (Guangyun, in Baxter's transcription), which glosses the first 'to explain' and 'to untie', and the second 'to let go' and 'to untie'.
\end{cousinweb}

\begin{grammarlens}[title={What you've built}]
One more word for saying what I do.
\end{grammarlens}

\subsection*{Your turn}
\begin{itemize}
\raggedright
  \item \emph{Say it:} \emph{jiěshì}
  \item \emph{Say it:} \emph{jiěshì}, once more
  \item \emph{Say it:} \emph{zhège dá'àn nǐ néng jiěshì ma} — Can you explain this answer?
  \item \emph{Recall:} say the Mandarin for to pass an exam, then the Mandarin for the answer to a question, then say \emph{jiěshì} again
\end{itemize}

\begin{tcolorbox}[breakable,colback=teal!4,colframe=teal!35!black,title={Before you move on}]
What does jiěshì mean? (\textquotedblleft{}To explain\textquotedblright{}.) Say it once more.
\end{tcolorbox}
